% ===================================================================
% Abschnitt "Abstrakte Formulierung für das Vorgehen"
% ===================================================================

\newglossaryentry{bem_ziel_statistischer_test}{
    name={Was ist das Ziel eines statistischen Tests?},
    description={Anhand von Daten zu entscheiden, ob eine zuvor formulierte Annahme über die Verteilung bestätigt werden kann.},
    type=dinge
}

\newglossaryentry{bem_bestandteile_testtheorie_vokabular}{
    name={Welche Bestandteile umfasst die abstrakte Formulierung eines Testproblems (neben dem statistischen Modell selbst)?},
    description={Hypothese ($\theta\in H_0$), Alternative ($\theta\in H_A$, mit $H_0\cap H_A=\varnothing$), Teststatistik, Test, Entscheidungsregel.},
    type=dinge
}

% ===================================================================
% Abschnitt "Wahl der Hypothesen"
% ===================================================================

\newglossaryentry{bem_gleiche_h0_verschiedene_ha}{
    name={Was zeigt das Mensa-Beispiel bezüglich der Wahl von $H_0$ und $H_A$ durch unterschiedliche Interessengruppen?},
    description={Beide Gruppen (Mensabetreiber, Studierendenvertreter) wählen dieselbe Nullhypothese $H_0:\theta=\theta_0$, aber unterschiedliche, jeweils zu ihrem Interesse passende einseitige Alternativen ($\theta<\theta_0$ bzw.\ $\theta>\theta_0$).},
    type=dinge
}

% ===================================================================
% Abschnitt "Neyman--Pearson-Formulierung der Testtheorie"
% ===================================================================

\newglossaryentry{def_testproblem_neyman_pearson}{
    name={Wie lautet die Neyman--Pearson-Formulierung eines Testproblems?},
    description={Gegeben ein statistisches Modell $\{P_\theta:\theta\in\Theta\}$ und $H_0,H_A\subseteq\Theta$ mit $H_0\cap H_A=\varnothing$: Hypothese $\theta\in H_0$, Alternative $\theta\in H_A$. Das Entscheidungsproblem "$H_0$ gegen $H_A$" heißt Testproblem, $H_0$ heißt Nullhypothese.},
    type=dinge
}

\newglossaryentry{bem_muss_h0_und_ha_ganz_theta_abdecken}{
    name={Müssen $H_0$ und $H_A$ zusammen den gesamten Parameterraum $\Theta$ abdecken?},
    description={Nein, es wird nur $H_0\cap H_A=\varnothing$ gefordert; $H_0\cup H_A=\Theta$ ist nicht notwendig.},
    type=dinge
}

\newglossaryentry{def_einfache_zusammengesetzte_hypothese}{
    name={Wann heißt eine Hypothese (bzw.\ Alternative) einfach, wann zusammengesetzt?},
    description={Ist $H_0$ (bzw.\ $H_A$) einelementig, heißt sie einfach; andernfalls zusammengesetzt.},
    type=dinge
}

\newglossaryentry{bem_wahl_nullhypothese_konservativ}{
    name={Wonach richtet sich die Wahl der Nullhypothese, und was bedeutet ihre Ablehnung?},
    description={Die Nullhypothese wird konservativ gewählt -- i.\,d.\,R.\ als das Gegenteil des vermuteten Zusammenhangs, den man belegen möchte. Wird $H_0$ im Test abgelehnt, gilt der vermutete Zusammenhang als statistisch signifikant.},
    type=dinge
}

% ===================================================================
% Abschnitt "Teststatistik" und "$p$-Wert"
% ===================================================================

\newglossaryentry{def_teststatistik_th}{
    name={Wie ist eine Teststatistik definiert, und nach welchem Kriterium wählt man sie?},
    description={Eine Teststatistik ist eine Abbildung $T:\mathcal X\to\mathbb R$. Man wählt sie so, dass große Werte $T(x)$ dafür sprechen, dass eine unter $H_0$ atypische Stichprobe vorliegt.},
    type=dinge
}

\newglossaryentry{bem_wahl_teststatistik_haengt_von_alternative_ab}{
    name={Warum kann dieselbe naheliegende Größe (z.\,B.\ $\bar X$) je nach Richtung der Alternative eine geeignete oder ungeeignete Teststatistik sein?},
    description={Ob große Werte einer Größe für die Alternative sprechen, hängt von der Richtung ab, in die die Alternative vom Nullhypothesen-Wert abweicht. Bei einer Alternative "kleiner als $\theta_0$" ist z.\,B.\ nicht $\bar X$, sondern $1-\bar X$ die geeignete Teststatistik; bei einer Alternative "größer als $\theta_0$" ist dagegen $\bar X$ direkt geeignet.},
    type=dinge
}

\newglossaryentry{def_p_wert_th}{
    name={Wie ist der $p$-Wert einer Stichprobe $x$ definiert?},
    description={$p(x) = \sup_{\theta\in H_0} P_\theta\big(T(X)\ge T(x)\big)$.},
    type=dinge
}

\newglossaryentry{bem_p_wert_intuition_th}{
    name={Was drückt der $p$-Wert intuitiv aus?},
    description={Wie wahrscheinlich es unter der Nullhypothese maximal ist, eine mindestens so extreme Teststatistik zu beobachten wie die tatsächlich gemessene. Ein kleiner $p$-Wert spricht für die Alternative.},
    type=dinge
}

% ===================================================================
% Abschnitt "Test, Signifikanzniveau, Verwerfungsbereich"
% ===================================================================

\newglossaryentry{def_test_abbildung}{
    name={Wie ist ein (statistischer) Test formal definiert?},
    description={Eine Abbildung $\varphi:\mathcal X\to\{0,1\}$. $\varphi(x)=1$ bedeutet: $H_0$ wird verworfen; $\varphi(x)=0$ bedeutet: $H_0$ wird nicht verworfen.},
    type=dinge
}

\newglossaryentry{def_signifikanzniveau_th}{
    name={Wie ist das Signifikanzniveau eines Tests definiert?},
    description={$\sup_{\theta\in H_0} P_\theta(\varphi(X)=1)$ -- die maximale Wahrscheinlichkeit, dass der Test unter $H_0$ verwirft.},
    type=dinge
}

\newglossaryentry{bem_p_wert_mehr_information_als_niveau}{
    name={Warum gibt gute Statistik den $p$-Wert an statt nur das Testergebnis zu einem festen Niveau?},
    description={Der $p$-Wert enthält mehr Information als das bloße Signifikanzniveau -- er zeigt, wie deutlich (bzw.\ knapp) das Ergebnis ausfällt, nicht nur ob verworfen wird oder nicht.},
    type=dinge
}

\newglossaryentry{def_verwerfungsbereich_th}{
    name={Was ist der Verwerfungsbereich $C_\alpha$ eines Tests?},
    description={$C_\alpha=\{x\in\mathcal X:\varphi(x)=1\}$ -- die Menge aller Stichproben, für die der Test die Nullhypothese verwirft. Statt $\varphi$ auszuwerten, kann man äquivalent prüfen, ob $x\in C_\alpha$ liegt.},
    type=dinge
}

% ===================================================================
% Abschnitt "Fehler 1.\ und 2.\ Art"
% ===================================================================

\newglossaryentry{def_fehler_1_art}{
    name={Was ist ein Fehler 1.\ Art?},
    description={Eine gültige Hypothese wird verworfen: $\theta\in H_0$ und $\varphi(X)=1$.},
    type=dinge
}

\newglossaryentry{def_fehler_2_art}{
    name={Was ist ein Fehler 2.\ Art?},
    description={Eine ungültige Hypothese wird nicht verworfen: $\theta\in H_A$ und $\varphi(X)=0$.},
    type=dinge
}

\newglossaryentry{bem_fehler_1_art_beschraenkt_durch_niveau}{
    name={Wodurch ist die Wahrscheinlichkeit für einen Fehler 1.\ Art beschränkt?},
    description={Durch das Signifikanzniveau des Tests.},
    type=dinge
}

\newglossaryentry{bem_ziel_testkonstruktion}{
    name={Was ist das Ziel bei der Konstruktion eines Tests zu vorgegebenem Signifikanzniveau $\alpha$?},
    description={Einen Test zu finden, dessen Signifikanzniveau den vorgegebenen Wert $\alpha$ nicht überschreitet -- idealerweise, ohne dass die Wahrscheinlichkeit für einen Fehler 2.\ Art zu groß wird.},
    type=dinge
}

\newglossaryentry{bem_tradeoff_fehler_1_2_art}{
    name={Welcher Zielkonflikt (Trade-off) besteht zwischen Fehler 1.\ und Fehler 2.\ Art?},
    description={Die Wahrscheinlichkeit für einen Fehler 2.\ Art lässt sich reduzieren, aber nur um den Preis, dass die Wahrscheinlichkeit für einen Fehler 1.\ Art (bzw.\ das Signifikanzniveau) größer wird.},
    type=dinge
}

% ===================================================================
% Abschnitt "Vorgehen bei einem Hypothesentest" und "Asymmetrie"
% ===================================================================

\newglossaryentry{bem_rezept_hypothesentest}{
    name={Welche sieben Schritte umfasst das allgemeine Vorgehen bei einem Hypothesentest?},
    description={(1) statistisches Modell/Verteilungsannahme, (2) Hypothese und Alternative formulieren, (3) Signifikanzniveau wählen, (4) Teststatistik konstruieren, (5) Entscheidungsregel festlegen, (6) Test auf die Stichprobe anwenden, (7) Schlussfolgerung und $p$-Wert angeben.},
    type=dinge
}

\newglossaryentry{bem_asymmetrie_testentscheidung}{
    name={Warum bedeutet "$H_0$ nicht verwerfen" keine Aussage über die Alternative?},
    description={Weil eine Entscheidung FÜR eine Hypothese grundsätzlich weniger sicher wäre als eine Entscheidung GEGEN eine Hypothese. "Nicht verwerfen" heißt nur, dass die Stichprobe nicht ausreicht, um $H_0$ mit hinreichender Sicherheit abzulehnen -- nicht, dass $H_0$ dadurch bewiesen wäre.},
    type=dinge
}

% ===================================================================
% Abschnitt "Begleitbeispiel: Geburtenstatistik" und Ein-/zweiseitige Tests
% ===================================================================

\newglossaryentry{bem_konstruktion_teststatistik_ggz_zgws}{
    name={Wie motiviert man die Konstruktion einer Teststatistik für einen Test über eine Erfolgswahrscheinlichkeit $\theta$ mit Hilfe von Gesetz der großen Zahlen und zentralem Grenzwertsatz?},
    description={Unter $H_0$ sollte $\left|\frac1n S_n-\theta_0\right|$ klein sein (Gesetz der großen Zahlen: $\frac1n S_n\to\theta_0$). Zusätzlich ist $\sqrt{\frac{n}{\theta_0(1-\theta_0)}}\left(\frac1n S_n-\theta_0\right)$ für großes $n$ näherungsweise standardnormalverteilt (zentraler Grenzwertsatz). Man wählt den Betrag dieser standardisierten Größe als Teststatistik.},
    type=dinge
}

\newglossaryentry{bem_geburtenbeispiel_typ}{
    name={Welcher Testtyp wird durchgeführt, wenn $H_0:\theta=\theta_0$ gegen $H_A:\theta\ne\theta_0$ getestet wird (wie im Geburtenstatistik-Beispiel)?},
    description={Ein zweiseitiger Gauß-Test.},
    type=dinge
}

\newglossaryentry{bem_einseitig_vs_zweiseitig_unterscheidung}{
    name={Woran erkennt man, ob ein Test einseitig oder zweiseitig ist?},
    description={Bei einer Alternative der Form $\theta<\theta_0$ oder $\theta>\theta_0$ (wie im Mensa-Beispiel) ist der Test einseitig; bei einer Alternative $\theta\ne\theta_0$ (wie im Geburten-Beispiel) ist er zweiseitig.},
    type=dinge
}

% ===================================================================
% Abschnitt "Einfacher Gauß-Test"
% ===================================================================

\newglossaryentry{bem_gauss_test_voraussetzung}{
    name={Unter welcher Voraussetzung ist der einfache Gauß-Test anwendbar?},
    description={$X_1,\dots,X_n$ normalverteilt mit unter $H_0$ bekannter Varianz $\sigma^2$.},
    type=dinge
}

\newglossaryentry{bem_gauss_test_bei_bernoulli}{
    name={Wie lässt sich der Gauß-Test auch bei Bernoulli-verteilten Stichprobenvariablen anwenden?},
    description={Näherungsweise für großes $n$: Unter $H_0$ kennt man $\theta_0$ und damit auch die Varianz $\sigma^2=\theta_0(1-\theta_0)$.},
    type=dinge
}

\newglossaryentry{formel_teststatistik_gauss_test}{
    name={Wie lautet die Teststatistik des einfachen Gauß-Tests?},
    description={$T(x) = \sqrt n\,\dfrac{\bar x-\mu_0}{\sigma}$.},
    type=dinge
}

\newglossaryentry{bem_verteilung_teststatistik_gauss_normal_vs_bernoulli}{
    name={Welche Verteilung hat $T(X)$ beim Gauß-Test, je nachdem ob die Stichprobenvariablen normal- oder Bernoulli-verteilt sind?},
    description={Bei normalverteilten Stichprobenvariablen ist $T(X)$ exakt der Betrag einer standardnormalverteilten Zufallsvariable; bei Bernoulli-verteilten Stichprobenvariablen nur näherungsweise (für großes $n$).},
    type=dinge
}

\newglossaryentry{formel_gauss_test_zweiseitig}{
    name={Wie lautet die Entscheidungsregel des zweiseitigen Gauß-Tests zum Signifikanzniveau $\alpha$?},
    description={$\varphi(x)=1$ falls $T(x)\ge z_{1-\alpha/2}$, sonst $0$. Verwerfungsbereich: $|T|>z_{1-\alpha/2}$.},
    type=dinge
}

\newglossaryentry{formel_gauss_test_einseitig}{
    name={Wie lauten die Entscheidungsregeln der beiden einseitigen Varianten des Gauß-Tests?},
    description={Für $H_0:\theta\le\theta_0$ gegen $H_A:\theta>\theta_0$: verwerfe falls $T(x)\ge z_{1-\alpha}$. Für $H_0:\theta\ge\theta_0$ gegen $H_A:\theta<\theta_0$: verwerfe falls $T(x)\le -z_{1-\alpha}$.},
    type=dinge
}

% ===================================================================
% Abschnitt "Einfacher $t$-Test"
% ===================================================================

\newglossaryentry{bem_t_test_voraussetzung}{
    name={Unter welcher Voraussetzung ist der einfache $t$-Test (statt des Gauß-Tests) anwendbar?},
    description={$X_1,\dots,X_n$ normalverteilt mit unbekannter Varianz $\sigma^2$.},
    type=dinge
}

\newglossaryentry{bem_t_test_vorgehen_analogie}{
    name={Wie unterscheidet sich das Vorgehen beim $t$-Test von dem beim Gauß-Test?},
    description={Es ist analog, nur wird die unbekannte Varianz in der Teststatistik durch ihren Schätzer $\hat\sigma^2=\frac{1}{n-1}\sum(x_i-\bar x)^2$ ersetzt, und man arbeitet mit den Quantilen der $t$-Verteilung statt der Standardnormalverteilung.},
    type=dinge
}

\newglossaryentry{formel_teststatistik_t_test}{
    name={Wie lautet die Teststatistik des einfachen $t$-Tests, und welcher Verteilung folgt sie unter $H_0$?},
    description={$T(x)=\sqrt n\,\dfrac{\bar x-\mu_0}{\hat\sigma}$; $T(X)$ ist der Betrag einer $t$-verteilten Zufallsvariable mit $n-1$ Freiheitsgraden.},
    type=dinge
}

\newglossaryentry{formel_t_test_entscheidungsregeln}{
    name={Wie unterscheiden sich die Entscheidungsregeln des $t$-Tests von denen des Gauß-Tests?},
    description={Man ersetzt $\sigma$ durch $\hat\sigma$ und die Normalquantile $z_{1-\alpha}$ bzw.\ $z_{1-\alpha/2}$ durch die entsprechenden $t$-Quantile $t_{n-1,1-\alpha}$ bzw.\ $t_{n-1,1-\alpha/2}$ mit $n-1$ Freiheitsgraden; die Struktur der Entscheidungsregeln bleibt sonst gleich.},
    type=dinge
}

% ===================================================================
% Abschnitt "Zwei-Stichprobentests"
% ===================================================================

\newglossaryentry{bem_zwei_stichproben_situation}{
    name={In welcher Situation kommen Zwei-Stichprobentests mit dem Differenzen-Trick zum Einsatz?},
    description={Wenn man zwei gepaarte Messungen $(X_i,Y_i)$ an derselben Versuchseinheit hat (z.\,B.\ zwei Messungen an derselben Person) und wissen möchte, ob sich die beiden Messgrößen systematisch unterscheiden. Dabei dürfen $X_i$ und $Y_i$ selbst abhängig sein -- nur die Paare $(X_i,Y_i)$ müssen über die Einheiten hinweg unabhängig sein.},
    type=dinge
}

\newglossaryentry{bem_differenzentrick}{
    name={Welcher Trick erlaubt es, den gepaarten Zwei-Stichprobentest auf einen gewöhnlichen Einstichproben-Test zurückzuführen?},
    description={Man bildet die Differenzen $D_i:=X_i-Y_i$. Da die Paare $(X_i,Y_i)$ über $i$ unabhängig sind, sind auch $D_1,\dots,D_n$ unabhängig (u.i.v.\ unter der Annahme gemeinsamer Normalverteilung mit EW und Varianz) -- man kann also den gewöhnlichen $t$-Test auf die $D_i$ anwenden, mit Nullhypothese $\mu_D=0$.},
    type=dinge
}

\newglossaryentry{formel_gepaarter_t_test_zweiseitig}{
    name={Wie lautet die Teststatistik und Entscheidungsregel des zweiseitigen gepaarten $t$-Tests?},
    description={Mit $\bar D=\frac1n\sum(x_i-y_i)$ und $\hat\sigma^2=\frac{1}{n-1}\sum(d_i-\bar d)^2$: verwerfe $H_0:\mu=0$ falls $\left|\sqrt n\,\dfrac{\bar D}{\hat\sigma}\right|\ge t_{n-1,1-\alpha/2}$.},
    type=dinge
}
